% 中英双语论文模板 —— 中文正文 + 英文题名/摘要/关键词（国内学报常用）
% 引擎：xelatex
\documentclass[11pt,a4paper,UTF8,fontset=fandol]{ctexart}

\usepackage{amsmath,amssymb,amsthm,mathtools}
% ctexart 已内置 xeCJK，无需重复加载；fontset=fandol 保证跨平台中文字体可用
\usepackage{geometry}
\usepackage{hyperref}
\geometry{margin=2.6cm}
\hypersetup{colorlinks=true,linkcolor=blue}
\setlength{\parindent}{2em}

\newtheorem{theorem}{定理}[section]
\newtheorem{lemma}[theorem]{引理}
\newcommand{\R}{\mathbb{R}}
\newcommand{\Om}{\Omega}

% 中文题录
\title{一类抛物型方程解的长时间行为}
\author{杨建伟$^{1}$}
\date{\today}

\begin{document}
\maketitle

\begin{abstract}
\noindent\textbf{摘要：}本文研究有界区域上带非线性源项的抛物型方程
\begin{equation}\label{eq:p}
  u_t-\Delta u+f(u)=g(x), \qquad (x,t)\in\Om\times(0,\infty),
\end{equation}
在 Dirichlet 边界条件下解的长时间行为。通过构造 Lyapunov 泛函并
建立一致的 $H^2$ 先验估计，证明了解半群存在有限维全局吸引子。

\smallskip
\noindent\textbf{关键词：}抛物型方程；全局吸引子；Lyapunov 泛函；先验估计

\smallskip
\noindent\textbf{MSC(2020)：}35K55, 35B41, 37L30
\end{abstract}

% ---------- 英文题录 ----------
\begin{center}
{\large\bfseries Long-Time Behavior of Solutions to a Class of
Parabolic Equations}\\[6pt]
{\normalsize Jianwei Yang$^{1}$}\\[3pt]
{\small $^{1}$School of Mathematics and Statistics, NCWU, Zhengzhou 450046, China}
\end{center}

\noindent\textbf{Abstract.} We study the long-time behavior of solutions
to the semilinear parabolic equation
$u_t-\Delta u+f(u)=g(x)$ in a bounded domain $\Om\subset\R^n$ with
Dirichlet boundary conditions. By constructing a Lyapunov functional
and establishing uniform $H^2$ a priori estimates, we prove the
existence of a finite-dimensional global attractor for the associated
solution semigroup. Our result extends earlier work on polynomially
growing nonlinearities to the critical exponent case.

\smallskip
\noindent\textbf{Keywords:} parabolic equation; global attractor;
Lyapunov functional; a priori estimate

\bigskip
\section{Introduction / 引言}

Consider the initial--boundary value problem
\begin{equation}\label{eq:ibvp}
  \begin{cases}
    u_t-\Delta u+f(u)=g(x), & (x,t)\in\Om\times\R^+,\\
    u=0, & (x,t)\in\partial\Om\times\R^+,\\
    u(x,0)=u_0(x), & x\in\Om ,
  \end{cases}
\end{equation}
where $\Om\subset\R^n$ is a bounded smooth domain and
$f\in C^1(\R)$ satisfies $f'(s)\ge -c_0$ and
$|f(s)|\le C(1+|s|^{p-1})$.

设 $\Om\subset\R^n$ 为有界光滑区域。方程 \eqref{eq:ibvp} 描述了
反应扩散过程中的浓度演化，其长时间动力学是无穷维动力系统理论中的
经典问题。

\begin{theorem}[Main Theorem / 主要定理]
Assume $g\in L^2(\Om)$ and $u_0\in H^1_0(\Om)$. Then the semigroup
$S(t)$ generated by \eqref{eq:ibvp} possesses a global attractor
$\mathcal{A}\subset H^2(\Om)\cap H^1_0(\Om)$ with finite fractal
dimension.
\end{theorem}

\begin{proof}[Proof / 证明]
The energy identity gives a bounded absorbing set in $H^1_0(\Om)$;
a second estimate obtained by testing with $-\Delta u$ yields
compactness in $H^1_0$. Finite dimensionality follows from the
squeezing property. \par
由能量等式可得 $H^1_0(\Om)$ 中的有界吸收集；再以 $-\Delta u$
为检验函数得到 $H^2$ 估计，从而由 Sobolev 紧嵌入获得渐近紧性。
挤压性质给出吸引子的有限维性。
\end{proof}

\section{Conclusion / 结论}
We have established the existence and finite dimensionality of the
global attractor. 本文建立了一类抛物型方程全局吸引子的存在性与
有限维性，后续工作将考虑非自治情形。

\end{document}
